% ---
% id: EX-CH01-002
% titre: Crible d'Ératosthène et formule d'Euler
% chapitre_id: 1
% chapitre_nom: Nombres entiers & divisibilité
% annee_scolaire: 9H
% difficulte: 2
% tags: [nombres-premiers, divisibilite]
% auteur: Bussy D.
% ---

\begin{exercice}{EX-CH01-002}{Crible d'Ératosthène et formule d'Euler}

\begin{mdframed}[backgroundcolor=vlmgray, innerleftmargin=10pt, innerrightmargin=10pt, innertopmargin=6pt, innerbottommargin=6pt, skipabove=-5pt, skipbelow=5pt]
Un \textbf{nombre premier} est un nombre entier supérieur à 1 qui ne possède que deux diviseurs : 1 et lui-même.
\end{mdframed}


\medskip
{\footnotesize Des diverses méthodes utilisées pour fabriquer une table des nombres premiers, la plus intéressante est celle que l'on attribue à Ératosthène, astronome, mathématicien et géographe grec (III\textsuperscript{e} siècle av.\ J.-C.). Cette méthode, qui a reçu le nom de \textit{crible d'Ératosthène}, consiste à écrire la suite des nombres entiers naturels à partir de 2 et à barrer dans cette suite les multiples de 2 à partir de 4, puis les multiples de 3 à partir de 6, ainsi de suite avec les multiples de 5, 7, 11, etc. Finalement, sont premiers les nombres qui restent non barrés. \hfill Gérard Charrière, \textit{L'algèbre mode d'emploi}}

\vfill
\textbf{A.}

\medskip
Applique la méthode du crible d'Ératosthène pour dresser la liste des nombres premiers inférieurs à 144.

\vspace{20pt}
\sol{\def\arraystretch{1.5}
\begin{tabularx}{\linewidth}{|*{12}{>{\centering\arraybackslash}X|}}
\hline
& 2 & 3 & 4 & 5 & 6 & 7  & 8 & 9 & 10 & 11  & 12  \\
\hline
13  & 14  & 15  & 16  & 17  & 18  & 19  & 20  & 21  & 22  & 23  & 24  \\
\hline
25  & 26  & 27  & 28  & 29  & 30  & 31  & 32  & 33  & 34  & 35  & 36  \\
\hline
37  & 38  & 39  & 40  & 41  & 42  & 43  & 44  & 45  & 46  & 47  & 48  \\
\hline
49  & 50  & 51  & 52  & 53  & 54  & 55  & 56  & 57  & 58  & 59  & 60  \\
\hline
61  & 62  & 63  & 64  & 65  & 66  & 67  & 68  & 69  & 70  & 71  & 72  \\
\hline
73  & 74  & 75  & 76  & 77  & 78  & 79  & 80  & 81  & 82  & 83  & 84  \\
\hline
85  & 86  & 87  & 88  & 89  & 90  & 91  & 92  & 93  & 94  & 95  & 96  \\
\hline
97  & 98  & 99  & 100 & 101 & 102 & 103 & 104 & 105 & 106 & 107 & 108 \\
\hline
109 & 110 & 111 & 112 & 113 & 114 & 115 & 116 & 117 & 118 & 119 & 120 \\
\hline
121 & 122 & 123 & 124 & 125 & 126 & 127 & 128 & 129 & 130 & 131 & 132 \\
\hline
133 & 134 & 135 & 136 & 137 & 138 & 139 & 140 & 141 & 142 & 143 & 144 \\
\hline
\end{tabularx}\def\arraystretch{1}}%
{\color{black}\def\arraystretch{1.5}\begin{tabularx}{\linewidth}{|*{12}{>{\centering\arraybackslash}X|}}
\hline
& 2 & 3 & $\Colxcancel[vlmgreen]{4}$ & 5 & $\Colxcancel[vlmgreen]{6}$ & 7 & $\Colxcancel[vlmgreen]{8}$ & $\Colxcancel[vlmgreen]{9}$ & $\Colxcancel[vlmgreen]{10}$ & 11 & $\Colxcancel[vlmgreen]{12}$ \\
\hline
13 & $\Colxcancel[vlmgreen]{14}$ & 15 & $\Colxcancel[vlmgreen]{16}$ & 17 & $\Colxcancel[vlmgreen]{18}$ & 19 & $\Colxcancel[vlmgreen]{20}$ & $\Colxcancel[vlmgreen]{21}$ & $\Colxcancel[vlmgreen]{22}$ & 23 & $\Colxcancel[vlmgreen]{24}$ \\
\hline
$\Colxcancel[vlmgreen]{25}$ & $\Colxcancel[vlmgreen]{26}$ & $\Colxcancel[vlmgreen]{27}$ & $\Colxcancel[vlmgreen]{28}$ & 29 & $\Colxcancel[vlmgreen]{30}$ & 31 & $\Colxcancel[vlmgreen]{32}$ & $\Colxcancel[vlmgreen]{33}$ & $\Colxcancel[vlmgreen]{34}$ & $\Colxcancel[vlmgreen]{35}$ & $\Colxcancel[vlmgreen]{36}$ \\
\hline
37 & $\Colxcancel[vlmgreen]{38}$ & $\Colxcancel[vlmgreen]{39}$ & $\Colxcancel[vlmgreen]{40}$ & 41 & $\Colxcancel[vlmgreen]{42}$ & 43 & $\Colxcancel[vlmgreen]{44}$ & $\Colxcancel[vlmgreen]{45}$ & $\Colxcancel[vlmgreen]{46}$ & 47 & $\Colxcancel[vlmgreen]{48}$ \\
\hline
$\Colxcancel[vlmgreen]{49}$ & $\Colxcancel[vlmgreen]{50}$ & $\Colxcancel[vlmgreen]{51}$ & $\Colxcancel[vlmgreen]{52}$ & 53 & $\Colxcancel[vlmgreen]{54}$ & $\Colxcancel[vlmgreen]{55}$ & $\Colxcancel[vlmgreen]{56}$ & $\Colxcancel[vlmgreen]{57}$ & $\Colxcancel[vlmgreen]{58}$ & 59 & $\Colxcancel[vlmgreen]{60}$ \\
\hline
61 & $\Colxcancel[vlmgreen]{62}$ & $\Colxcancel[vlmgreen]{63}$ & $\Colxcancel[vlmgreen]{64}$ & $\Colxcancel[vlmgreen]{65}$ & $\Colxcancel[vlmgreen]{66}$ & 67 & $\Colxcancel[vlmgreen]{68}$ & $\Colxcancel[vlmgreen]{69}$ & $\Colxcancel[vlmgreen]{70}$ & 71 & $\Colxcancel[vlmgreen]{72}$ \\
\hline
73 & $\Colxcancel[vlmgreen]{74}$ & $\Colxcancel[vlmgreen]{75}$ & $\Colxcancel[vlmgreen]{76}$ & $\Colxcancel[vlmgreen]{77}$ & $\Colxcancel[vlmgreen]{78}$ & 79 & $\Colxcancel[vlmgreen]{80}$ & $\Colxcancel[vlmgreen]{81}$ & $\Colxcancel[vlmgreen]{82}$ & 83 & $\Colxcancel[vlmgreen]{84}$ \\
\hline
$\Colxcancel[vlmgreen]{85}$ & $\Colxcancel[vlmgreen]{86}$ & $\Colxcancel[vlmgreen]{87}$ & $\Colxcancel[vlmgreen]{88}$ & 89 & $\Colxcancel[vlmgreen]{90}$ & $\Colxcancel[vlmgreen]{91}$ & $\Colxcancel[vlmgreen]{92}$ & $\Colxcancel[vlmgreen]{93}$ & $\Colxcancel[vlmgreen]{94}$ & $\Colxcancel[vlmgreen]{95}$ & $\Colxcancel[vlmgreen]{96}$ \\
\hline
97 & $\Colxcancel[vlmgreen]{98}$ & $\Colxcancel[vlmgreen]{99}$ & $\Colxcancel[vlmgreen]{100}$ & 101 & $\Colxcancel[vlmgreen]{102}$ & 103 & $\Colxcancel[vlmgreen]{104}$ & $\Colxcancel[vlmgreen]{105}$ & $\Colxcancel[vlmgreen]{106}$ & 107 & $\Colxcancel[vlmgreen]{108}$ \\
\hline
109 & $\Colxcancel[vlmgreen]{110}$ & $\Colxcancel[vlmgreen]{111}$ & $\Colxcancel[vlmgreen]{112}$ & 113 & $\Colxcancel[vlmgreen]{114}$ & $\Colxcancel[vlmgreen]{115}$ & $\Colxcancel[vlmgreen]{116}$ & $\Colxcancel[vlmgreen]{117}$ & $\Colxcancel[vlmgreen]{118}$ & $\Colxcancel[vlmgreen]{119}$ & $\Colxcancel[vlmgreen]{120}$ \\
\hline
$\Colxcancel[vlmgreen]{121}$ & $\Colxcancel[vlmgreen]{122}$ & $\Colxcancel[vlmgreen]{123}$ & $\Colxcancel[vlmgreen]{124}$ & $\Colxcancel[vlmgreen]{125}$ & $\Colxcancel[vlmgreen]{126}$ & 127 & $\Colxcancel[vlmgreen]{128}$ & $\Colxcancel[vlmgreen]{129}$ & $\Colxcancel[vlmgreen]{130}$ & 131 & $\Colxcancel[vlmgreen]{132}$ \\
\hline
133 & $\Colxcancel[vlmgreen]{134}$ & $\Colxcancel[vlmgreen]{135}$ & $\Colxcancel[vlmgreen]{136}$ & 137 & $\Colxcancel[vlmgreen]{138}$ & 139 & $\Colxcancel[vlmgreen]{140}$ & $\Colxcancel[vlmgreen]{141}$ & $\Colxcancel[vlmgreen]{142}$ & $\Colxcancel[vlmgreen]{143}$ & $\Colxcancel[vlmgreen]{144}$ \\
\hline
\end{tabularx}\def\arraystretch{1}}

\vspace{20pt}
Nombres premiers inférieurs à 144 : \sol{\dotfill

\dotfill

\dotfill}{2, 3, 5, 7, 11, 13, 17, 19, 23, 29, 31, 37, 41, 43, 47, 53, 59, 61, 67, 71, 73, 79, 83, 89, 97, 101, 103, 107, 109, 113, 127, 131, 137, 139}

\vspace{20pt}
{\footnotesize \og Jusqu'à ce jour, les mathématiciens ont en vain tenté de découvrir un ordre dans la suite des nombres premiers et nous avons des raisons de croire que c'est un mystère que l'esprit ne pénétrera jamais. \fg \hfill Leonhard Euler}

\medskip
{\footnotesize \og Euler rencontre en 1772 une curieuse expression algébrique : $n^2 + n + 41$. En remplaçant $n$ par des nombres entiers de 0 à 39, il obtient une suite ininterrompue de 40 nombres premiers. Mais cette suite se brise lorsqu'il remplace $n$ par 40. En effet, $40^2 + 40 + 41 = 1681$ n'est pas un nombre premier. \fg \hfill Gérard Charrière, \textit{L'algèbre mode d'emploi}}

\newpage
\textbf{B.}

\medskip
Retrouve la suite des 40 nombres premiers obtenus par Euler en remplaçant $n$ par les nombres entiers de 0 à 39 dans l'expression $n^2 + n + 41$.

\vspace{20pt}
\def\arraystretch{1.5}
\begin{tabularx}{\linewidth}{|*{10}{>{\centering\arraybackslash}X|}}
\hline
\cellcolor{vlmgray} 0 &\cellcolor{vlmgray} 1 &\cellcolor{vlmgray} 2 &\cellcolor{vlmgray} 3 &\cellcolor{vlmgray} 4 &\cellcolor{vlmgray} 5 &\cellcolor{vlmgray} 6 &\cellcolor{vlmgray} 7 &\cellcolor{vlmgray} 8 &\cellcolor{vlmgray}9 \\
\hline
\textcolor{\colSolution}{41} & \textcolor{\colSolution}{43} & \textcolor{\colSolution}{47} & \textcolor{\colSolution}{53} & \textcolor{\colSolution}{61} & \textcolor{\colSolution}{71} & \textcolor{\colSolution}{83} & \textcolor{\colSolution}{97} & \textcolor{\colSolution}{113} & \textcolor{\colSolution}{131} \\
\hline
\multicolumn{10}{X}{}\\
\hline
\cellcolor{vlmgray} 10 &\cellcolor{vlmgray} 11 &\cellcolor{vlmgray} 12 &\cellcolor{vlmgray} 13 &\cellcolor{vlmgray} 14 &\cellcolor{vlmgray} 15 &\cellcolor{vlmgray} 16 &\cellcolor{vlmgray} 17 &\cellcolor{vlmgray} 18 &\cellcolor{vlmgray}19 \\
\hline
\textcolor{\colSolution}{151} & \textcolor{\colSolution}{173} & \textcolor{\colSolution}{197} & \textcolor{\colSolution}{223} & \textcolor{\colSolution}{251} & \textcolor{\colSolution}{281} & \textcolor{\colSolution}{313} & \textcolor{\colSolution}{347} & \textcolor{\colSolution}{383} & \textcolor{\colSolution}{421} \\
\hline
\multicolumn{10}{X}{}\\
\hline
\cellcolor{vlmgray} 20 &\cellcolor{vlmgray} 21 &\cellcolor{vlmgray} 22 &\cellcolor{vlmgray} 23 &\cellcolor{vlmgray} 24 &\cellcolor{vlmgray} 25 &\cellcolor{vlmgray} 26 &\cellcolor{vlmgray} 27 &\cellcolor{vlmgray} 28 &\cellcolor{vlmgray}29 \\
\hline
\textcolor{\colSolution}{461} & \textcolor{\colSolution}{503} & \textcolor{\colSolution}{547} & \textcolor{\colSolution}{593} & \textcolor{\colSolution}{641} & \textcolor{\colSolution}{691} & \textcolor{\colSolution}{753} & \textcolor{\colSolution}{797} & \textcolor{\colSolution}{853} & \textcolor{\colSolution}{911} \\
\hline
\multicolumn{10}{X}{}\\
\hline
\cellcolor{vlmgray} 30 &\cellcolor{vlmgray} 31 &\cellcolor{vlmgray} 32 &\cellcolor{vlmgray} 33 &\cellcolor{vlmgray} 34 &\cellcolor{vlmgray} 35 &\cellcolor{vlmgray} 36 &\cellcolor{vlmgray} 37 &\cellcolor{vlmgray} 38 &\cellcolor{vlmgray}39 \\
\hline
\textcolor{\colSolution}{971} & \textcolor{\colSolution}{1033} & \textcolor{\colSolution}{1097} & \textcolor{\colSolution}{1163} & \textcolor{\colSolution}{1231} & \textcolor{\colSolution}{1301} & \textcolor{\colSolution}{1373} & \textcolor{\colSolution}{1447} & \textcolor{\colSolution}{1523} & \textcolor{\colSolution}{1601} \\
\hline
\end{tabularx}\def\arraystretch{1}

\sol{}{\vspace{20pt}Pour $n = 40$ : $40^2 + 40 + 41 = 1600 + 40 + 41 = 1681 = 41^2$, qui n'est pas premier.}
\end{exercice}